\chapter*{Glossary}\addcontentsline{toc}{chapter}{Glossary}
\markboth{Glossary}{Glossary}
These definitions belong to the scope used in this book. A term can have a different technical meaning in another discipline; the model and context must travel with it.

\begin{description}
\item[Action.] A specified physical transformation, with its quantities, boundary and conditions. A proposed action is not evidence of an executed event.
\item[Actor.] The person, organisation or explicitly declared decision rule choosing an action. A value calculation does not itself make that choice.
\item[Affected support.] Every potential factor whose value can change in the action. It can extend beyond the immediately touched stocks in a coupled model.
\item[Attraction.] Approach towards a specified set or state under a specified dynamic law. A minimum of a function alone supplies no attraction.
\item[Attribution.] A declared assignment of a joint account to component actions. Mathematical closure does not establish causality, fairness or ownership.
\item[Baseline.] The state used for an action comparison. In a driven event, the frozen action baseline follows any declared prior external change.
\item[Boundedness.] Remaining within finite bounds. In this book's fixed-mass domain, stock and potential bounds already follow from the physical constraints.
\item[Capacity, experimental.] A prospective account balance limiting affordability under a separately declared rule. It is distinct from physical storage or throughput capacity.
\item[Conservation.] Constancy of a specified total within a declared boundary. A conserved total can be distributed in a functionally unsuitable way.
\item[Curvature.] How a slope changes as its input changes. In the quadratic potential it supplies an exact finite correction to the initial marginal estimate.
\item[Derivative.] The limiting rate of change of a function with respect to one input. Its unit is the function's unit divided by the input's unit.
\item[Dissipation.] Irreversible physical degradation of available driving differences. Its representation and valuation require a declared physical model; it is not a synonym for every monetary expense.
\item[Drive.] A declared external change to the represented physical state. Its potential effect must be distinguished from the actor's subsequent action value.
\item[EBU.] The project's name for a proposed account of physical action. Its core finite expression is an affected potential difference minus a separately justified process burden. The letters do not establish an energy unit or a currency.
\item[Equilibrium.] A state left unchanged by a specified dynamic law. Thermodynamic equilibrium has the more specific physical meaning discussed in Chapter~\ref{ch:balance}.
\item[Externality.] An effect on others that is not fully reflected in the private costs or benefits guiding a decision.
\item[Feasibility.] Satisfaction of the declared physical constraints. It does not by itself establish permission, affordability or desirability.
\item[Field, directional.] The initial marginal reduction of potential per unit quantity in a specified action direction. It can change over a finite action.
\item[Functional balance.] Conditions compatible with specified functioning, potentially maintained through continuous throughput. It need not mean equal stocks or immobility.
\item[Gaussian ruler.] Here, a quadratic measure of deviation associated with the shape of a Gaussian reference density. It does not assert normally distributed physical trajectories.
\item[Gradient.] The vector of partial derivatives of a scalar function. In this separable model it is the vector of local marginal burdens.
\item[Hessian.] The matrix of second derivatives. For the fixed Gaussian ruler it is diagonal, with positive entries $1/\sigma_i^2$.
\item[Homeostasis.] Active regulation of selected internal variables within ranges compatible with function. It depends on mechanisms and resources.
\item[Identity.] An equality following from stated definitions and assumptions. Its exactness does not demonstrate the empirical adequacy of those assumptions.
\item[Local information.] Information available within a declared physical or informational neighbourhood. Complete affected support and limited actor knowledge are different requirements.
\item[Lossless transfer.] A transfer in which source decrease equals destination increase for the represented resource. Real losses require explicit physical accounting.
\item[Marginal burden.] The sensitivity of potential to one stock, holding other coordinates fixed. Its sign indicates the local slope, not a moral judgement.
\item[Normal distribution.] A probability distribution with a bell-shaped density. Choosing its geometry as a ruler is distinct from fitting it to observations.
\item[Permission.] An institutional or account rule governing which feasible actions may proceed. It must be stated separately from valuation.
\item[Potential.] The declared scalar function of physical state used for valuation. In this book it is dimensionless deviation, not automatically thermodynamic energy or social welfare.
\item[Process burden.] A separately declared nonnegative physical contribution in compatible account units, subject to the no-double-counting boundary. The ideal lossless illustration sets it to zero.
\item[Public good.] In the ideal economic definition, a benefit that is nonrival and nonexcludable. Actual services may have congestion and exclusion constraints.
\item[Recovery.] A return of specified variables or functions to a declared condition after disturbance, with a stated criterion and time horizon. It requires dynamics and evidence beyond an account identity.
\item[Reference.] The declared comparison state. Compatibility with mass and capacities does not establish reachability or equilibrium.
\item[Scale.] A fixed positive stock quantity setting the ruler's sensitivity to displacement. It is not automatically an observed standard deviation or a safety limit.
\item[Settlement.] Application of separately declared rules to account balances after an event. Field value, compensation and personal responsibility are distinct questions.
\item[Stability.] A specified form of controlled response to perturbation under a specified dynamic law. One must name the relevant definition; boundedness alone is insufficient.
\item[Steady state.] Constancy of represented stocks or macroscopic variables over time. It may coexist with continuing input and output flows.
\item[Telescoping.] Cancellation of intermediate endpoint terms in a sequence of differences evaluated under one fixed potential.
\item[Throughput.] Physical flow through a system over time. Constant stock does not imply zero throughput.
\end{description}

\chapter*{Symbols and units}\addcontentsline{toc}{chapter}{Symbols and units}
\markboth{Symbols and units}{Symbols and units}
The stock unit in the illustrations is the litre (L). A different homogeneous resource may use another stock unit. Account quantities below are dimensionless for the Gaussian construction. Superscript $\star$ labels a reference; $\mathsf T$ transposes a vector; superscript $n$ labels an event state. None of these means a square.

\renewcommand{\arraystretch}{1.16}
\begin{longtable}{@{}p{65pt}p{194pt}p{80pt}@{}}
\toprule Symbol & Meaning & Unit here\\\midrule\endfirsthead
\toprule Symbol & Meaning & Unit here\\\midrule\endhead
\bottomrule\endfoot
$i,j$ & Cell indices & Labels\\
$x_i,\ x$ & Stock at a cell; full stock vector & L\\
$z$ & Frozen post-drive action baseline & L\\
$x_i^\star$ & Declared reference stock & L\\
$\sigma_i>0$ & Fixed scale & L\\
$M$ & Conserved total stock & L\\
$K_i$ & Physical storage capacity & L\\
$a,\ e,\ G$ & Action, directed edge, action group & Labels/set\\
$q$ & Finite transferred quantity & L\\
$r,\ \Delta t$ & Transfer rate; duration & L/min; min\\
$e_i$ & Unit vector selecting cell $i$ & Dimensionless\\
$S_e=-e_i+e_j$ & Lossless transfer direction & Dimensionless\\
$\delta x_a$ & Physical increment of action $a$ & L\\
$\delta x_G$ & Sum of declared group increments & L\\
$T_a(z)$ & State after action map $a$ & L\\
$A$ & Complete affected potential support & Set\\
$V_i,\ V,\ V_A$ & Local, total and affected potential & Account unit\\
$\mu_i$ & $\partial V/\partial x_i$: marginal burden & L$^{-1}$\\
$\mu=\nabla V$ & Vector of partial derivatives & L$^{-1}$\\
$f_e=\mu_i-\mu_j$ & Directional marginal field & L$^{-1}$\\
$H$ & Fixed Hessian, diagonal $1/\sigma_i^2$ & L$^{-2}$\\
$p_i(u)$ & Gaussian reference density & L$^{-1}$\\
$C_a$ & Separate residual process burden & Account unit\\
$E,\ \Delta E_a$ & Finite signed action account & Account unit\\
$E_G$ & Joint field reduction (zero cost) & Account unit\\
$E_a^{(0)}$ & Singleton value at common start & Account unit\\
$R_a$ & Declared common-path attribution & Account unit\\
$\lambda$ & Fraction along straight interpolation & Dimensionless\\
$n,\ N$ & Cell count (Ch.12); event index and transition count (Ch.17) & Counts\\
$D_{{\rm ext},n}$ & Signed external potential injection & Account unit\\
$B_i$ & Prospective experimental balance & Account unit\\
$J$ & Signed cumulative external audit & Account unit\\
\end{longtable}
\renewcommand{\arraystretch}{1}

In $u^{\mathsf T}v$, multiply matching vector entries and add them. In $\sum_i$, add over the declared indices. In $\int$, accumulate the specified continuously varying quantity over its stated interval. These operations are introduced where first needed; they do not select an actor policy.
